

% Dies definiert die Klasse Ihres Dokuments. Ersetzen Sie 'PHKAinf' mit dem tatsächlichen Namen Ihrer Dokumentklasse.
\documentclass{PHKAinf}

% Hier binden wir notwendige Pakete ein. Diese Zeile muss vor '\begin{document}' stehen.
% Es ist wichtig, dass die Datei 'packages_macros.tex' keine Befehle enthält, die erst nach '\begin{document}' ausgeführt werden dürfen.
\input{packages_macros}

% Zusätzliche Pakete können hier eingebunden werden, falls 'packages_macros.tex' nicht alle benötigten Pakete abdeckt.
\usepackage[round]{natbib}
\usepackage[ngerman]{babel}
\usepackage{listings}
\usepackage{algorithm2e}
\usepackage{xcolor}

% Hier beginnt der Inhalt Ihres Dokuments.
\begin{document}

\Titlepage %This generates the titlepage
\frontmatter
\input{Rahmenseiten/Danksagung}
\cleardoublepage
\input{Rahmenseiten/Abstracts}
\cleardoublepage

\tableofcontents
\listoftables
\listoffigures

\mainmatter

% In der Datei "IhreArbeit.tex" fügen Sie den Hauptinhalt Ihres Dokuments ein.
\input{IhreArbeit}

% Literaturverzeichnis
\bibliographystyle{apalike}
\bibliography{Literaturverzeichnis}

% Glossare und Abkürzungen
\printglossary[type=\acronymtype]
\printglossary

\appendix
% Anhänge
\input{Rahmenseiten/Anhang}

\end{document}
